\begin{eyenumerate}
\item \PubItem[\with M. Fornasier and G. Savaré]{Density of subalgebras of Lipschitz functions in metric Sobolev spaces and applications to Wasserstein Sobolev spaces}
  	 {\textit{Journal of Functional Analysis}  | 285.11 (2023)}
  	% {We introduce and investigate a notion of multivalued $\lambda$-dissipative probability vector field (MPVF) in the Wasserstein space $\mathcal{P}_2(\mathsf{X})$ of Borel probability measures on a Hilbert space $\mathsf{X}$. Taking inspiration from the theory of dissipative operators in Hilbert spaces and of Wasserstein gradient flows of geodesically convex functionals, we study local and global well posedness of evolution equations driven by dissipative MPVFs. Our approach is based on a measure-theoretic version of the Explicit Euler scheme, for which we prove novel convergence results with optimal error estimates under an abstract CFL stability condition, which do not rely on compactness arguments and also hold when $\mathsf{X}$ has infinite dimension. We characterize the limit solutions by a suitable Evolution Variational Inequality (EVI), inspired by the Bénilan notion of integral solutions to dissipative evolutions in Banach spaces. Existence, uniqueness and stability of EVI solutions are then obtained under quite general assumptions, leading to the generation of a semigroup of nonlinear contractions.}
\item \PubItem[]{The general class of Wasserstein Sobolev spaces: density of cylinder functions, reflexivity, uniform convexity and Clarkson's inequalities}
  	  {\textit{Calculus of Variations and Partial Differential Equations}  | 62.7 (2023)}
\item \PubItem[\with G. Cavagnari and G. Savaré]{Dissipative PVFs and generation of evolution semigroups in Wasserstein spaces}
  	 {\textit{Probability Theory and Related Fields}  | 185.3-4 (2023), pp. 1087-1182}
  	% {We introduce and investigate a notion of multivalued $\lambda$-dissipative probability vector field (MPVF) in the Wasserstein space $\mathcal{P}_2(\mathsf{X})$ of Borel probability measures on a Hilbert space $\mathsf{X}$. Taking inspiration from the theory of dissipative operators in Hilbert spaces and of Wasserstein gradient flows of geodesically convex functionals, we study local and global well posedness of evolution equations driven by dissipative MPVFs. Our approach is based on a measure-theoretic version of the Explicit Euler scheme, for which we prove novel convergence results with optimal error estimates under an abstract CFL stability condition, which do not rely on compactness arguments and also hold when $\mathsf{X}$ has infinite dimension. We characterize the limit solutions by a suitable Evolution Variational Inequality (EVI), inspired by the Bénilan notion of integral solutions to dissipative evolutions in Banach spaces. Existence, uniqueness and stability of EVI solutions are then obtained under quite general assumptions, leading to the generation of a semigroup of nonlinear contractions.}
 \item \PubItem[\with G. Savaré]{A simple relaxation approach to duality for OT problems in completely regular spaces}
   {\textit{Journal of Convex Analysis} 29.1 (2022), pp.1-12}
   %{In this paper, we present a simple and direct approach to duality for Optimal Transport for lower semicontinuous cost functionals in arbitrary completely regular topological spaces, showing that the Optimal Transport functional can be interpreted as the largest sublinear and weakly lower semicontinuous functional extending the cost between pairs of Dirac masses.}
   \item \PubItem[\with M. D'Amico, J. De Tullio and G. Osimo]{Mathematical Analysis - Module 1 Exercises}
   {Egea - Le dispense del Pellicano (2021)}
   %{First volume of the BAI series, a collection of exercises books for the Bachelor in ``Mathematical and Computing Sciences for Artificial Intelligence'' at Bocconi University.}
 \item \PubItem[\with M.Martini]{Numerical methods for a system of coupled Cahn-Hilliard equations}
   {\textit{Communications in Applied and Industrial Mathematics}, 12, (2021), Issue 1, pp. 1-12}
  % {In this work, we consider a system of coupled Cahn-Hilliard equations describing the phase separation of a copolymer and a homopolymer blend. We propose some numerical methods to approximate the solution of the system which are based on suitable combinations of existing schemes for the single Cahn-Hilliard equation. As a verification for our experimental approach, we present some tests and a detailed description of the numerical solutions’ behaviour obtained by varying the values of the system’s characteristic parameters.}
 
    %\item \PubItem[\with N. Surname, N. Surname]{Title}
   %{\arXiv{code} | \emph{Journal} 2019 | links}
   %{Description}
\end{eyenumerate}
